
% STABILITY SECTION FOR PAPER I

\subsection{Perturbative Stability}
\label{sec:stability}

A critical requirement for any dark energy model is perturbative stability.
For scalar field models, this is characterized by the sound speed squared:
%
\begin{equation}
c_s^2 = \frac{P_X}{P_X + 2X P_{XX}}
\end{equation}
%
where $P(X,\phi)$ is the Lagrangian density and $X = \frac{1}{2}(\partial\phi)^2$.

For our canonical kinetic term $\mathcal{L} = X - V(\phi)$, we have $P_X = 1$ 
and $P_{XX} = 0$, yielding:
%
\begin{equation}
c_s^2 = 1
\end{equation}
%
\textit{at all redshifts}, regardless of the potential $V(\phi)$.

This result is particularly important given the phantom-like behavior ($w < -1$) observed 
in the late-time evolution of our model (see Sec. \ref{sec:quintessence}). unlike standard phantom models 
which achieve $w < -1$ by flipping the kinetic sign (leading to ghosts and instabilities), 
the Evaporating Universe maintains a canonical kinetic term. The apparent phantom equation of state 
arises from the specific dynamics of the potential slope and energy exchange, not from a 
pathological kinetic sector.

Thus, the model guarantees:
\begin{itemize}
    \item \textbf{Gradient stability}: $c_s^2 = 1 > 0$
    \item \textbf{Causality}: $c_s^2 = 1$ (luminal propagation)
    \item \textbf{Ghost-free}: Positive kinetic energy density
\end{itemize}
This confirms that the Evaporating Universe is physically viable and free from the pathologies 
often associated with dynamic dark energy models.
